\documentclass[11pt,a4paper]{article}
\usepackage[margin=1in]{geometry}
\usepackage{hyperref}
\usepackage{enumitem}
\usepackage{graphicx}
\usepackage{xcolor}

% -----------------------------------------------------
% bvraghav-tiet-ucs503p-article.sty
% -----------------------------------------------------

% Variable: Lab Instructor
\makeatletter \newcommand{\BvrSetLabInstructor}[1]{%
  \gdef\Bvr@LabInstructor{#1}}
% default
\newcommand{\Bvr@LabInstructor}{Dr. Raghav %
  B. Venkataramaiyer}
% public accessor
\newcommand{\BvrLabInstructorValue}{\Bvr@LabInstructor}
\makeatother

% Add subtitle
\usepackage{titling}
% -- Set title bold
\pretitle{\begin{center}\bfseries\fontsize{18bp}{18bp}\selectfont}
\makeatletter
\newcommand{\BvrSubtitle}[1]{%
  \posttitle{%
    \par\end{center}
    \begin{center}\large#1\end{center}
    \vskip 0.5em}%
}
\makeatother

% Hints in the section title(s)
\newcommand{\BvrHint}[1]{%
  {\normalfont\normalsize\itshape\hfill #1}}

% -----------------------------------------------------

\title{Smart Library Seat Management System}

\BvrSetLabInstructor{Dr. Raghav B. Venkataramaiyer}

\BvrSubtitle{%
  Project Proposal for UCS503P  \\
  Submitted to: \BvrLabInstructorValue \\ \bfseries
  Thapar Institute of Engineering and Technology% 
}

\author{
  Student Name, CSED%
  \thanks{Roll No: \texttt{1024030xxx}, %
    Email: \texttt{student\_be24\@thapar.edu}}
  \and 
  Student Name 2, CSED%
  \thanks{Roll No: \texttt{1024030xxx}, %
    Email: \texttt{student2\_be24\@thapar.edu}}
}

\date{\today}

\begin{document}
\maketitle

\section{Higher-order goal%
  \BvrHint{\ldots secondary framing}}
This project aims to optimize the utilization of physical study spaces in educational institutions by providing real-time visibility and fair management of library seating. The immediate engineering objective is to translate \emph{efficient resource allocation and occupancy tracking} into a measurable, deployable software solution.

\section{Time-to-value %
\BvrHint{\ldots primary heuristic}}
\begin{itemize}[leftmargin=0.7cm]
    \item We will deliver meaningful increments quickly by prototyping the core user workflow (seat viewing and temporary-away status) and validating it with users within a short iteration window.
    \item We will prioritize fast feedback over perfection by using weekly iterations (and targeting CI-backed automation early) to reduce release friction.
    \item Success is defined by what can be delivered and measured in the first iteration, then improved based on user feedback.
\end{itemize}

\section{Problem Statement}
Students currently have to physically roam around the library to find an empty seat. A major pain point is that many students leave the library temporarily for lunch, calls, or washroom breaks, but keep their belongings on their seats for long periods, effectively blocking those seats for other students.
\begin{itemize}[leftmargin=0.7cm]
    \item \textbf{High Search Friction:} Students waste time navigating multiple floors to find available spots.
    \item \textbf{Unfair Resource Blocking:} Unattended belongings block seats indefinitely.
    \item \textbf{Low Transparency:} No real-time data on library occupancy levels.
    \item \textbf{Suboptimal Space Utilization:} Librarians lack analytics to understand peak usage and effectively manage the layout.
\end{itemize}

\textbf{Why this matters now (contextual relevance):} In an educational setting with limited seating bandwidth, intelligent monitoring and policy-enforced release of temporarily abandoned seats can noticeably improve fairness and reduce wasted capacity.

\section{Proposed Solution %
\BvrHint{\ldots Software Proposition}}
\subsection{Overview}
Build a \textbf{web-based} ``Library Seat Management'' system with the following components:
\begin{itemize}[leftmargin=0.7cm]
    \item \textbf{Real-Time Floor Map:} Interactive floor plan displaying live status (Available, Occupied, Temporarily Away).
    \item \textbf{Temporary Away Permit:} Allows students to lock their seat for a short configurable time with a countdown timer.
    \item \textbf{Smart Seat Finder:} AI-ranked recommendations based on preferences like quiet zones, power outlets, or group seating.
    \item \textbf{Fair-Use and No-Show Protection:} Automatic seat release if students fail to check-in or overstay their break limits.
    \item \textbf{Admin/Librarian Dashboard:} Live occupancy tracking, analytics, and policy configurations.
\end{itemize}

\subsection{Core workflow}
\begin{enumerate}[leftmargin=0.7cm]
    \item Student scans ID card at the entrance to initiate an active session.
    \item Student selects an available seat from the digital floor map.
    \item Student activates a 'Temporary Away' permit if they need to leave for a short break (e.g., 30 mins).
    \item System warns the user before expiration, and automatically releases the seat if the timer runs out.
    \item Student scans ID at the exit to conclude the session and free up the seat.
\end{enumerate}

\subsection{Operational constraints for an educational setting}
\begin{itemize}[leftmargin=0.7cm]
    \item Keep the solution usable on low-to-mid range mobile devices via responsive web design.
    \item Avoid dependence on expensive physical IoT sensors initially; focus on software-based workflows, with architecture supporting future IoT integration.
    \item Design for deployability using common web hosting and managed databases.
\end{itemize}

\section{Solution Approach %
\BvrHint{\ldots Engineering focus}}
This is an engineering course project, so we focus on scalable administrative and workflow aspects.

\subsection{Fairness and Abandoned Seat Detection %
\BvrHint{\ldots non-research}}
Detecting seat abandonment will be state-based:
\begin{itemize}[leftmargin=0.7cm]
    \item Track sessions through ID card entries/exits or digital check-ins.
    \item Use simple timers to monitor 'Temporary Away' statuses.
    \item Automatically transition seat status and alert admins based on simple configurable thresholds (e.g., >30 mins away).
\end{itemize}

\subsection{Web architecture %
\BvrHint{\ldots web-based solution}}
A typical 3-tier web architecture:
\begin{itemize}[leftmargin=0.7cm]
    \item \textbf{Frontend:} Responsive UI for students (floor maps, dashboard) and staff (analytics, overrides).
    \item \textbf{Backend API:} Authentication, session management, timer handling, seat allocation logic.
    \item \textbf{Database:} Store users, seats, historical sessions, and layout metadata.
\end{itemize}

\subsection{CI/CD and fast delivery enabler}
CI/CD will be used to:
\begin{itemize}[leftmargin=0.7cm]
    \item Automatically build and run tests on every change.
    \item Produce repeatable deployment artifacts to staging.
    \item Enable safe and frequent releases of incremental features.
\end{itemize}

\section{Evaluation Criterion %
\BvrHint{\ldots measurable and attributable}}
We define success using measurable metrics tied to the core workflow.

\subsection{Primary evaluation metric}
\textbf{Seat Search Time:} time between a user entering the library and checking into an available seat.
\begin{itemize}[leftmargin=0.7cm]
    \item Target: median search time $\le$ 3 minutes during pilot window.
    \item Attribution: measured from system check-in logs.
\end{itemize}

\subsection{Secondary evaluation metrics}
\begin{itemize}[leftmargin=0.7cm]
    \item \textbf{Seat Utilization Rate:} Percentage of time available seats are actively occupied vs blocked.
    \item \textbf{Abandoned Seat Recovery:} Number of seats automatically released due to expired temporary away permits.
    \item \textbf{User Satisfaction:} Student reported satisfaction with finding seats easily.
    \item \textbf{Reliability:} Uptime and availability of the real-time map during peak hours.
\end{itemize}

\subsection{Pilot validation plan %
\BvrHint{\ldots high-level}}
\begin{itemize}[leftmargin=0.7cm]
    \item Recruit pilot users to test the digital flow in a specific library zone.
    \item Compare search times and seat utilization metrics against baseline observations.
\end{itemize}

\section{Scalability %
\BvrHint{\ldots with foresight}}
The proposition should scale in both theoretical and practical senses.

\begin{enumerate}[leftmargin=0.7cm]
    \item \textbf{Scalable (theoretically):} The system should support high concurrency during peak exam hours:
    \begin{itemize}[leftmargin=0.7cm]
        \item Optimizing real-time status updates via WebSockets or efficient polling.
        \item Caching seat layout metadata.
        \item Using stateless API design.
    \end{itemize}
    
    \item \textbf{Operationally deployable (practically):} The system should run on common web hosting:
    \begin{itemize}[leftmargin=0.7cm]
        \item managed database,
        \item containerized or platform-hosted deployment,
        \item CI/CD pipeline for staging/production.
    \end{itemize}
\end{enumerate}

\section{Engine availability heuristic}
A mature set of ``engines'' (tooling/services) is available to implement this as a web-based project:
\begin{itemize}[leftmargin=0.7cm]
    \item Web framework for REST APIs and templated/reactive pages.
    \item Real-time engines (e.g., Socket.io, Firebase, Supabase Realtime).
    \item Managed relational database.
    \item Test framework for unit/integration tests.
    \item CI runner and staging deployment target.
\end{itemize}
We will use these mature components to focus on workflow design, correctness, and measurement.

\section{Project scope and deliverables %
\BvrHint{\ldots iteration-friendly}}
\subsection{Initial deliverable %
\BvrHint{\ldots first iteration}}
\begin{itemize}[leftmargin=0.7cm]
    \item User login and session creation.
    \item Basic interactive seat selection from a list or simple grid.
    \item The 'Temporary Away' timer workflow.
    \item CI: build + backend unit tests + linting.
\end{itemize}

\subsection{Subsequent deliverables}
\begin{itemize}[leftmargin=0.7cm]
    \item Interactive visual floor plan mapping.
    \item Smart seat recommendation engine.
    \item Waitlisting and group seating logic.
    \item Admin dashboard for real-time analytics.
\end{itemize}

\section{Risks and mitigations}
\begin{itemize}[leftmargin=0.7cm]
    \item \textbf{Adoption Friction:} Students may forget to check in/out. Mitigate by using a very short initial grace period or integrating with existing turnstile/ID scanners if possible.
    \item \textbf{Real-Time Sync Issues:} High concurrent traffic might delay updates. Mitigate by using robust real-time caching solutions.
    \item \textbf{Physical Verification:} Digital status might not match physical reality. Mitigate with community reporting (e.g., "This seat is empty but marked occupied").
\end{itemize}

\section{Summary}
This project proposes a deployable web platform to manage library seating efficiently. It meets the course criteria by emphasizing time-to-value through rapid iterations, implementing CI/CD to support frequent safe releases, and using measurable evaluation criteria (specifically reducing seat search time and increasing active utilization). It remains focused on engineering workflow and scalability readiness while leaving room for future IoT/sensor enhancements.

\end{document}
